\section{Build and Test Infrastructure, Quality Control, and Disclosure}
\label{sec:buildtest}

\paragraph{Build and test infrastructure.}
\sysname is developed in a public monorepo, and every change is submitted as a
pull request that must pass continuous integration before it is reviewed. The
suite includes unit tests for the tooling of each agent team, integration
tests that exercise capture, replication, and analysis end to end against a
small reference infrastructure, and regression tests derived from every
previously confirmed finding, so that a capability once demonstrated cannot
be silently lost. Because the outputs are produced by non-deterministic
agents, the integration suite executes each scenario multiple times and
asserts on \emph{distributions} of outcomes rather than on exact traces, and
the run-to-run stability is tracked as a reported metric. The reference
infrastructures are themselves infrastructure-as-code, versioned alongside the source, so any result can be reproduced from a
commit hash.

\paragraph{Quality control and security of new content.}
Contributions require the review of a project maintainer, and they must pass
automated checks before they are merged, including static analysis and
linting, dependency scanning against advisory databases, and a verification
that any new dependency is recorded in the SBOM of the project. Releases are
tagged, signed, and published together with a generated SBOM and a provenance
attestation, following the recommended practices of CISA and NSA~\cite{cisansa}
and targeting the best-practices criteria of the OpenSSF~\cite{openssf} and
the build-provenance levels of SLSA~\cite{slsa}. Content generated by agents receives specific scrutiny. For example, twin
configurations, mocks, and synthetic datasets each carry a manifest that
records what produced them and with what confidence, and the synthetic data
is validated by the constraint agents of Section~\ref{sec:thrust2} and
checked to contain no material derived from production records. Twins are
held in isolated networks, access to them is restricted to the project team
and to the owning organization, and they are destroyed at the end of an
engagement.

\paragraph{Licensing and upstream contribution.}
\sysname will be released under a permissive OSI-approved license matching
those of the anchor products (BSD-2-Clause for nginx, Apache~2.0 for ZMap),
so contributions can flow upstream without license friction, and the reference infrastructures and their ground-truth knowledge
graphs will be released under a permissive data license. Improvements to the
anchor products are developed as upstream contributions from the start,
following the process of each project rather than being maintained as a fork.
For example, the ZMap extensions will be submitted as pull requests with
tests~\cite{zmapsrc}, and the nginx findings will follow the nginx security
process. Where a finding concerns a deployment configuration rather than
component code, the deliverable is detection tooling and documentation
contributed to the community.

\paragraph{Responsible disclosure and dual use.}
\sysname automates the discovery of exploitable chains in widely deployed
open-source software, and we treat the resulting dual-use considerations
explicitly. All findings follow coordinated vulnerability disclosure, in that a confirmed
upstream defect is reported privately to the security contact of that
project, together with a reproducer and a proposed remediation, and
publication is withheld until a fix is available or until a standard embargo
period has elapsed. Findings specific to the deployment of a partner belong
to that partner. The assessment of any system is performed only under a written authorization
from its owner, specifying the scope, the permitted observation techniques,
and the rules of engagement. The campus evaluations of Section~\ref{sec:evaluation} are conducted under
agreement with the responsible IT organizations at UCSB and ASU, with the
impact interlocks of Section~\ref{sec:thrust1} enforced throughout, and we
will seek IRB review where any activity touches human-subjects data. We will
publish the framework, the reference infrastructures, and the methods in
full, because defenders cannot use what they cannot obtain, and because the
capability we automate is one that sophisticated adversaries already
possess~\cite{measuring}. However, we will not publish weaponized exploits
for unpatched defects, and the released artifacts require an authorization
scope before an assessment can be run.
